\documentclass[11pt,a4paper]{article}
\usepackage{fontspec,amsmath,amssymb,amsthm,microtype}
\setmainfont{lmroman10-regular.otf}[BoldFont=lmroman10-bold.otf,ItalicFont=lmroman10-italic.otf,BoldItalicFont=lmroman10-bolditalic.otf]
\setmonofont{lmmono10-regular.otf}[Scale=MatchLowercase]
\usepackage[margin=25mm]{geometry}
\usepackage[unicode,hidelinks]{hyperref}
\setlength{\emergencystretch}{3em}
\setlength{\parskip}{0.4em}
\setlength{\parindent}{0pt}
\allowdisplaybreaks[2]
\raggedbottom
\title{\texorpdfstring{Actual KIPLD position laws and particle-\allowbreak{}independent contraction}{Actual KIPLD position laws and particle-independent contraction}}
\author{LangevinLean mathematical exposition}
\date{}
\begin{document}\maketitle
\noindent Source module: \texttt{KIPLMC/\allowbreak{}Flow\allowbreak{}Law.\allowbreak{}lean}. The accompanying HTML pairs this mathematical exposition with the exact Lean statements and records the scope of the formalization.


The scaled configuration space is \(\mathbb R^{d_\theta+Nd_x}\), with potential \(\bar U_N=V_N/N\). Its gradient is the actual physical-time particle force. The original pair potential's strong convexity and gradient Lipschitz constants pass to \(\bar U_N\) unchanged. Consequently the shared kinetic flow theorem applies with the original \(m,L,\gamma\), independently of the number of particles. This module constructs a probability law by an independent product and pushforward, and proves contraction of those actual laws. It does not assume Gibbs stationarity or identify an arbitrary forcing law with Brownian motion.


\section{\texorpdfstring{Use the Borel structure of continuous forcing paths}{Use the Borel structure of continuous forcing paths}}


The velocity input lives in the compact-path space \texttt{Velocity\allowbreak{}Forcing T}. Its natural measurable structure is its Borel structure. The two local instances install that structure and its defining compatibility witness. This makes the product of an initial velocity with a continuous forcing path an ordinary measurable source space. These are typeclass instances, not probabilistic independence hypotheses; independence will come from an explicit product measure.



{\raggedright\noindent Lean declarations: \texttt{forcing\allowbreak{}Measurable}, \texttt{forcing\allowbreak{}Borel}.\par}

\section{\texorpdfstring{Construct the physical-\allowbreak{}time position law}{Construct the physical-time position law}}


For an initial position law \(\mu\) and a joint source law \(Q\) of velocity and continuous input, form \(\mu\otimes Q\). Evolve each pair with the selected kinetic integral flow for \(\nabla\bar U_N\) and project to its position at time \(t\). The resulting pushforward is \texttt{kipld\allowbreak{}Position\allowbreak{}Law}. Its definition calls the shared \texttt{kinetic\allowbreak{}Position\allowbreak{}Law}, supplying the proved Lipschitz gradient of \(\bar U_N\). The normalization of the forcing subtracts its value at zero, so every continuous source path gives the correct initial condition. The Brownian specialization will choose the physical amplitude \(\sqrt{2\gamma/N}\); the definition deliberately admits any source law.



For \(0\le t\le T\), the endpoint is measurable because the selected flow depends continuously on its initial state and forcing. The product of probability inputs is a probability measure, and a measurable pushforward preserves total mass. \texttt{kipld\allowbreak{}Position\allowbreak{}Law\_\allowbreak{}probability} invokes exactly this proved shared result.



{\raggedright\noindent Lean declarations: \texttt{kipld\allowbreak{}Position\allowbreak{}Law}, \texttt{kipld\allowbreak{}Position\allowbreak{}Law\_\allowbreak{}probability}.\par}

\section{\texorpdfstring{Transfer the true Wasserstein infimum without a new stochastic axiom}{Transfer the true Wasserstein infimum without a new stochastic axiom}}


Let \(\mu,\nu\) have finite \(W_2\), and use the same source distribution \(Q\), independent of position for each evolution. If \(\gamma>0\) and \(\gamma^2\ge m+L\), then

\[W_2(P_t^Q\mu,P_t^Q\nu)\le\sqrt2e^{-mt/\gamma}W_2(\mu,\nu).\]

Lean applies the shared actual-law theorem to \(\bar U_N\). The three analytic inputs are proved earlier: differentiability, strong convexity with parameter \(m\), and smoothness with constant \(L\). The theorem extends every initial coupling by the common independent source, proves its output marginals, bounds its cost using the actual integral-equation contraction, and takes the infimum. Thus the result concerns the actual pushforward laws. Finiteness of the output transport cost follows from the bound and the explicit finite initial-cost hypothesis before conversion from extended to real distance. To replace the second output by the Gibbs target requires the separate stationarity theorem; this statement does not hide that obligation.



{\raggedright\noindent Lean declarations: \texttt{kipld\allowbreak{}Position\allowbreak{}Law\_\allowbreak{}wasserstein2\allowbreak{}Real}.\par}

\end{document}
